% preamble.tex — Shared preamble for STAT 212 Lecture Notes

% === Core Packages ===
\usepackage[utf8]{inputenc}
\usepackage[T1]{fontenc}
\usepackage{lmodern}
\usepackage[margin=1in]{geometry}

% === Math Packages ===
\usepackage{amsmath, amssymb, amsthm}
\usepackage{mathtools}
\usepackage{mathrsfs}
\usepackage{bm}
\usepackage{bbm}

% === Formatting ===
\usepackage{enumitem}
\usepackage{hyperref}
\usepackage{cleveref}
\usepackage{xcolor}
\usepackage[most]{tcolorbox}
\usepackage{tikz}
\usepackage{fancyhdr}

% === Hyperref Setup ===
\hypersetup{
    colorlinks=true,
    linkcolor=blue!60!black,
    urlcolor=blue!60!black,
    citecolor=blue!60!black,
}

% === Math Commands ===
\newcommand{\E}{\mathbb{E}}
\newcommand{\Prob}{\mathbb{P}}
\newcommand{\R}{\mathbb{R}}
\newcommand{\N}{\mathbb{N}}
\newcommand{\Z}{\mathbb{Z}}
\newcommand{\Q}{\mathbb{Q}}
\newcommand{\T}{\mathbb{T}}
\newcommand{\F}{\mathcal{F}}
\newcommand{\G}{\mathcal{G}}
\newcommand{\B}{\mathcal{B}}
\newcommand{\cH}{\mathcal{H}}
\newcommand{\cS}{\mathcal{S}}
\DeclareMathOperator{\Var}{Var}
\DeclareMathOperator{\Cov}{Cov}
\newcommand{\ind}{\mathbbm{1}}
\newcommand{\dto}{\xrightarrow{d}}
\newcommand{\pto}{\xrightarrow{p}}
\newcommand{\asto}{\xrightarrow{\text{a.s.}}}
\newcommand{\Lpto}[1]{\xrightarrow{L^{#1}}}
\newcommand{\abs}[1]{\left|#1\right|}
\newcommand{\norm}[1]{\left\|#1\right\|}
\newcommand{\dmu}{\,d\mu}
\newcommand{\dP}{\,d\Prob}

% === Color Definitions ===
% Definition: Blue
\definecolor{defbg}{HTML}{E8F0FE}
\definecolor{defframe}{HTML}{1A73E8}

% Theorem: Orange
\definecolor{thmbg}{HTML}{FFF3E0}
\definecolor{thmframe}{HTML}{E65100}

% Lemma: Teal
\definecolor{lembg}{HTML}{E0F2F1}
\definecolor{lemframe}{HTML}{00796B}

% Corollary: Purple
\definecolor{corbg}{HTML}{F3E5F5}
\definecolor{corframe}{HTML}{7B1FA2}

% Remark: Gray
\definecolor{rembg}{HTML}{F5F5F5}
\definecolor{remframe}{HTML}{616161}

% Example: Green
\definecolor{exbg}{HTML}{E8F5E9}
\definecolor{exframe}{HTML}{2E7D32}

% Claim: Indigo
\definecolor{clmbg}{HTML}{E8EAF6}
\definecolor{clmframe}{HTML}{3949AB}

% === tcolorbox Environments ===
\tcbuselibrary{theorems, breakable, skins}

% Definition
\newtcbtheorem[number within=chapter]{definition}{Definition}{
    enhanced, breakable,
    colback=defbg, colframe=defframe, colbacktitle=defframe, coltitle=white,
    fonttitle=\bfseries, boxrule=1pt, arc=2pt,
    left=8pt, right=8pt, top=6pt, bottom=6pt,
    attach boxed title to top left={yshift=-2mm, xshift=4mm},
    boxed title style={boxrule=0.5pt, arc=1pt},
    separator sign={.},
}{def}

% Theorem
\newtcbtheorem[use counter from=definition]{theorem}{Theorem}{
    enhanced, breakable,
    colback=thmbg, colframe=thmframe, colbacktitle=thmframe, coltitle=white,
    fonttitle=\bfseries, boxrule=1pt, arc=2pt,
    left=8pt, right=8pt, top=6pt, bottom=6pt,
    attach boxed title to top left={yshift=-2mm, xshift=4mm},
    boxed title style={boxrule=0.5pt, arc=1pt},
    separator sign={.},
}{thm}

% Lemma
\newtcbtheorem[use counter from=definition]{lemma}{Lemma}{
    enhanced, breakable,
    colback=lembg, colframe=lemframe, colbacktitle=lemframe, coltitle=white,
    fonttitle=\bfseries, boxrule=1pt, arc=2pt,
    left=8pt, right=8pt, top=6pt, bottom=6pt,
    attach boxed title to top left={yshift=-2mm, xshift=4mm},
    boxed title style={boxrule=0.5pt, arc=1pt},
    separator sign={.},
}{lem}

% Corollary
\newtcbtheorem[use counter from=definition]{corollary}{Corollary}{
    enhanced, breakable,
    colback=corbg, colframe=corframe, colbacktitle=corframe, coltitle=white,
    fonttitle=\bfseries, boxrule=1pt, arc=2pt,
    left=8pt, right=8pt, top=6pt, bottom=6pt,
    attach boxed title to top left={yshift=-2mm, xshift=4mm},
    boxed title style={boxrule=0.5pt, arc=1pt},
    separator sign={.},
}{cor}

% Remark
\newtcbtheorem[use counter from=definition]{remark}{Remark}{
    enhanced, breakable,
    colback=rembg, colframe=remframe, colbacktitle=remframe, coltitle=white,
    fonttitle=\bfseries, boxrule=0.5pt, arc=2pt,
    left=8pt, right=8pt, top=6pt, bottom=6pt,
    attach boxed title to top left={yshift=-2mm, xshift=4mm},
    boxed title style={boxrule=0.5pt, arc=1pt},
    separator sign={.},
}{rem}

% Example
\newtcbtheorem[use counter from=definition]{example}{Example}{
    enhanced, breakable,
    colback=exbg, colframe=exframe, colbacktitle=exframe, coltitle=white,
    fonttitle=\bfseries, boxrule=0.5pt, arc=2pt,
    left=8pt, right=8pt, top=6pt, bottom=6pt,
    attach boxed title to top left={yshift=-2mm, xshift=4mm},
    boxed title style={boxrule=0.5pt, arc=1pt},
    separator sign={.},
}{ex}

% Claim
\newtcbtheorem[use counter from=definition]{claim}{Claim}{
    enhanced, breakable,
    colback=clmbg, colframe=clmframe, colbacktitle=clmframe, coltitle=white,
    fonttitle=\bfseries, boxrule=0.5pt, arc=2pt,
    left=8pt, right=8pt, top=6pt, bottom=6pt,
    attach boxed title to top left={yshift=-2mm, xshift=4mm},
    boxed title style={boxrule=0.5pt, arc=1pt},
    separator sign={.},
}{clm}

% Proof styling — gray left border
\tcolorboxenvironment{proof}{
    blanker, breakable,
    left=5mm,
    before skip=10pt, after skip=10pt,
    borderline west={1pt}{0pt}{gray!60},
}

% === cleveref support for tcolorbox ===
\crefname{tcb@cnt@definition}{Definition}{Definitions}
\crefname{tcb@cnt@theorem}{Theorem}{Theorems}
\crefname{tcb@cnt@lemma}{Lemma}{Lemmas}
\crefname{tcb@cnt@corollary}{Corollary}{Corollaries}
\crefname{tcb@cnt@remark}{Remark}{Remarks}
\crefname{tcb@cnt@example}{Example}{Examples}
\crefname{tcb@cnt@claim}{Claim}{Claims}
